% from ega3-1-fr.tex
\oldpage[0\textsubscript{III}]{18}
\begin{env}[10.1.2]
\label{0.10.1.2-fr}
Supposons vérifiée l'une des hypothèses suivantes~:
\begin{enumerate}
  \item[(i)] $\mathfrak{J}$ est nilpotent~;
  \item[(ii)] $A$ est noethérien, $\mathfrak{J}$ est contenu dans le radical de
    $A$, et $M$ est de type fini.
\end{enumerate}
Alors les propriétés suivantes sont équivalentes~:
\begin{enumerate}
  \item[$a)$] $M$ est un $A$-module libre.
  \item[$b)$] $M/\mathfrak{J}M=M\otimes_A(A/\mathfrak{J})$ est un
    $(A/\mathfrak{J})$-module libre, et
    $\operatorname{Tor}_1^A(M,A/\mathfrak{J})=0$.
  \item[$c)$] $M/\mathfrak{J}M$ est un $(A/\mathfrak{J})$-module libre et
    l'homomorphisme canonique (\hyperref[0.10.1.1.2-fr]{10.1.1.2}) est
    injectif (donc bijectif).
\end{enumerate}
\end{env}

\begin{env}[10.1.3]
\label{0.10.1.3-fr}
Supposons que $A/\mathfrak{J}$ soit un \emph{corps} (autrement dit que
$\mathfrak{J}$ soit maximal), et que l'une des hypothèses (i), (ii) de
(\hyperref[0.10.1.2-fr]{10.1.2}) soit vérifiée. Alors les propriétés suivantes
sont équivalentes~:
\begin{enumerate}
  \item[$a)$] $M$ est un $A$-module libre.
  \item[$b)$] $M$ est un $A$-module projectif.
  \item[$c)$] $M$ est un $A$-module plat.
  \item[$d)$] $\operatorname{Tor}_1^A(M,A/\mathfrak{J})=0$.
  \item[$e)$] L'homomorphisme canonique
    (\hyperref[0.10.1.1.2-fr]{10.1.1.2}) est bijectif.
\end{enumerate}
Ce résultat s'appliquera en particulier dans les deux cas suivants~:
\begin{enumerate}
  \item[(i)] $M$ est un module \emph{quelconque} sur un anneau local $A$ dont
    l'idéal maximal $\mathfrak{J}$ est \emph{nilpotent} (par exemple un anneau
    local artinien).
  \item[(ii)] $M$ est un module \emph{de type fini} sur un anneau local
    \emph{noethérien}.
\end{enumerate}
\end{env}

\subsection{Critères locaux de platitude}
\label{subsection:0.10.2-fr}

\begin{env}[10.2.1]
\label{0.10.2.1-fr}
Les hypothèses et notations étant celles de
(\hyperref[0.10.1.1-fr]{10.1.1}), considérons les conditions suivantes~:
\begin{enumerate}
  \item[$a)$] $M$ est un $A$-module plat.
  \item[$b)$] $M/\mathfrak{J}M$ est un $(A/\mathfrak{J})$-module plat et
    $\operatorname{Tor}_1^A(M,A/\mathfrak{J})=0$.
  \item[$c)$] $M/\mathfrak{J}M$ est un $(A/\mathfrak{J})$-module plat et
    l'homomorphisme canonique (\hyperref[0.10.1.1.2-fr]{10.1.1.2}) est
    bijectif.
  \item[$d)$] Pour tout $n\geq1$, $M/\mathfrak{J}^nM$ est un
    $(A/\mathfrak{J}^n)$-module plat.
\end{enumerate}
On a alors les implications
\[
  a\Rightarrow b\Rightarrow c\Rightarrow d
\]
et si $\mathfrak{J}$ est \emph{nilpotent}, les quatre conditions $a)$, $b)$,
$c)$, $d)$ sont \emph{équivalentes}. Il en est de même si $A$ est noethérien
et en outre si $M$ est \emph{idéalement séparé}, c'est-à-dire que pour tout
idéal $\mathfrak{a}$ de $A$, le $A$-module $\mathfrak{a}\otimes_A M$ est
\emph{séparé} pour la topologie $\mathfrak{J}$-pré-adique.
\end{env}

\begin{env}[10.2.2]
\label{0.10.2.2-fr}
Soient $A$ un anneau noethérien, $B$ une $A$-algèbre commutative noethérienne,
$\mathfrak{J}$ un idéal de $A$ tel que $\mathfrak{J}B$ soit contenu dans le
radical de $B$, $M$ un $B$-module de type fini. Alors, lorsque $M$ est
considéré comme $A$-module, les conditions $a)$, $b)$, $c)$, $d)$ de
(\hyperref[0.10.2.1-fr]{10.2.1}) sont équivalentes.

Ce résultat s'applique surtout lorsque $A$ et $B$ sont des anneaux
\emph{locaux} noethériens, l'homomorphisme $A\to B$ un homomorphisme
\emph{local}. Plus particulièrement, si $\mathfrak{J}$ est alors l'idéal
\emph{maximal} de $A$, on peut, dans les conditions $b)$ et $c)$, supprimer
l'hypothèse que $M/\mathfrak{J}M$ est plat, qui est automatiquement vérifiée,
et la condition $d)$ signifie que les modules $M/\mathfrak{J}^nM$ sont
\emph{libres} sur les $A/\mathfrak{J}^n$.
\end{env}

